% !TEX program = xelatex
% Статья для сборника НИТУ «МИСиС» (шаблон: sadekovlar/misis-article-latex)
% Компиляция: xelatex article.tex (дважды) или tectonic article.tex

\documentclass{misisarticle}

% --- Классификационные коды (как в таблице курса) ---
\grnti{28.23, 06.73.75}
\asjc{1702, 1706}
\oecd{1.02}

% --- Название ---
\title{Классификация потока кредитных клиентов на основе ансамбля нейросетевых и классических методов}

% --- Автор ---
\misisauthor{И.~А.~Абрамов}{кафедра инженерной кибернетики НИТУ «МИСиС»}{Москва, Россия}{m2210595@edu.misis.ru}

% --- Аннотация и ключевые слова (зачтённые) ---
\abstract{%
  Задача — классификации кредитных клиентов на основе ML инструментов. Минимизируем ошибки первого и второго рода. Способ создания ансамбля моделей (или менеджера) учитывающего сильные и слабые стороны подходов foundation-модель TabPFN-2.5 и методы градиентного бустинга XGBoost, LightGBM, CatBoost. 
}

\keywords{%
  табличные foundation-модели, кредитный скоринг, дисбаланс классов, ансамбли моделей, ошибки классификации
}

\begin{document}

\maketitle

% ============================================================
\section{Введение}
\label{sec:intro}

Кредитный скоринг является задачей бинарной классификации с неравноценными ошибками. Выданный дефолтному клиенту кредит означает потерю выданных средств (ошибка I рода), отказ платёжеспособному клиенту~--- недополученную прибыль (ошибка II рода). Дефолты при этом редки: в используемом нами наборе их 3,55\%. Выбор модели, снижающей ошибку I рода при приемлемой ошибке II рода, напрямую определяет потери банка.

Стандартом для табличных данных остаются градиентные бустинги над решающими деревьями: XGBoost \cite{xgboost}, LightGBM \cite{lightgbm} и CatBoost \cite{catboost}. Недавно появилась альтернатива~--- табличная foundation-модель TabPFN, которая не обучается на целевых данных: обучающая выборка подаётся как контекст, а ответ вычисляется одним прямым проходом сети; на малых выборках она превосходит бустинги \cite{hollmann2025}. Для кредитного риска такие модели сравнивались с классическими в \cite{baesens2026}, где они использовались без настройки параметров и показали конкурентоспособное качество. Для кредитного скоринга предлагались и ансамбли, например стекинг на основе нейросетевой модели TabNet \cite{wang2024}.
% TODO: сверить формулировку про \cite{baesens2026} с полным текстом препринта (по аннотации: преимущество foundation-моделей заметнее на малых выборках и при дисбалансе классов).

Однако эти работы сравнивают модели по набору задач «в среднем». На практике качество зависит от режима, который задают объём выборки, число признаков и дисбаланс классов, а у каждой модели есть свои слабые места: бустинги уступают foundation-модели на малых выборках, а сама foundation-модель пропускает через сеть всю обучающую выборку при каждом выводе, и её стоимость растёт вместе с объёмом данных. Ни одна модель не может быть лучшей во всех режимах. В известных нам работах правило выбора между foundation-моделью и бустингом по режиму данных, выраженное через ошибки I и II рода, не предложено; пробел и сравнение с литературой подробно разобраны в разделе~\ref{sec:related}.

Цель работы~--- построить ансамблевый «менеджер моделей», который по характеристикам данных (объём выборки, число признаков, дисбаланс классов) отдаёт предпочтение TabPFN-2.5 либо одному из бустингов, и обосновать методику подбора его порогов по ошибкам I и II рода. Проверяется гипотеза: при заданной ошибке II рода менеджер даёт меньшую ошибку I рода, чем лучшая одиночная модель, усреднение моделей и стекинг.

Вклад работы:
\begin{itemize}
  \item профили ошибок I и II рода и ROC AUC четырёх моделей в зависимости от объёма выборки, доли дефолтов и числа признаков и слабые места моделей, выведенные из анализа ошибок (разд.~\ref{sec:results});
  \item методика подбора порогов переключения менеджера перебором по ошибкам I и II рода, подтверждённая графиками, на которых виден минимум (разд.~\ref{sec:method}, \ref{sec:results});
  \item сравнение менеджера с лучшей одиночной моделью, усреднением и стекингом на тесте, отложенном по времени (разд.~\ref{sec:results});
  \item воспроизводимый пайплайн: агрегация кредитных историй, разбиение по времени, подвыборки и код экспериментов \cite{repo}.
\end{itemize}

Далее в статье: разд.~\ref{sec:related}~--- обзор литературы, разд.~\ref{sec:method}~--- метод, разд.~\ref{sec:exp}~--- данные и протокол экспериментов, разд.~\ref{sec:results}~--- результаты, разд.~\ref{sec:concl}~--- выводы и ограничения.

% ============================================================
\section{Обзор литературы}
\label{sec:related}
% TODO: сгруппировать работы по темам, не <<статья за статьёй>>; закончить пробелом (gap) и мостом к вашему вкладу.

% ============================================================
\section{Метод}
\label{sec:method}
% TODO: критерий оптимальности, менеджер моделей, методика расчёта порогов переключения.

% ============================================================
\section{Эксперименты}
\label{sec:exp}
\subsection{Данные}

Обучающая часть набора содержит 3\,000\,000 заявок в 12 файлах, всего около 25~млн строк, в среднем 8,3 строки на заявку. Строка описывает один прошлый кредит клиента, метки заявок хранятся отдельно. Дефолтом считается просрочка 90 и более дней в течение 12 месяцев после выдачи; доля дефолтов в наборе равна 3,55\%. Анкетных данных (возраст, доход) в наборе нет, только кредитная история. Период сбора данных не раскрывается.

Исходные записи агрегированы до одной строки на заявку, получено 128 признаков, из них 43 категориальных. Для каждой заявки считаются: число и доля кредитов с каждым из флагов; число кредитов с каждым значением категориальных кодов кредита и число различных значений; число каждого статуса платежа в трёх окнах по номерам столбцов платёжной истории (0--2, 3--11, 12--24); первое, последнее (в порядке кредитов в истории) и число различных значений интервальных признаков; три последних статуса платежа по самому позднему кредиту. Интервальные признаки закодированы случайной нумерацией интервалов, поэтому порядок номеров не несёт смысла: для них не считаются среднее, максимум и минимум, а сами признаки трактуются как категориальные. Три столбца статусов платежей закодированы со сдвигом на единицу, и сдвиг устранён при агрегации. Смысл кодов статусов платежей в описании набора не раскрыт и в работе не интерпретируется.

Идентификаторы заявок упорядочены по времени, поэтому разбиение выполнено по времени (табл.~\ref{tab:split}): первые 60\% заявок образуют обучающий пул, следующие 20\%~--- валидацию, последние 20\%~--- тест. Валидация и тест представляют собой случайные подвыборки по 50\,000 заявок из своих временных окон. Доля дефолтов снижается от обучающего пула к тесту, то есть между частями есть временной дрейф. Кроме того, часть значений справочников (в частности, два значения типа кредита) встречается только в поздних файлах набора.

\begin{table}[t]
\caption{Разбиение данных по времени}
\label{tab:split}
\centering
\begin{tabular}{lrr}
\toprule
Часть & Объектов & Доля дефолтов, \% \\
\midrule
Обучающий пул & 1\,800\,000 & 3,61 \\
Валидация & 50\,000 & 3,49 \\
Тест & 50\,000 & 3,21 \\
\bottomrule
\end{tabular}
\end{table}

Обучающие выборки формируются случайными подвыборками из обучающего пула в трёх измерениях: размер $N \in \{2\,000;\ 5\,000;\ 10\,000;\ 25\,000;\ 50\,000\}$, доля дефолтов в выборке (1\%, естественная 3,61\%, 10\%, 25\%) и число признаков $k \in \{20, 50, 100, 128\}$; наборы признаков вложены друг в друга и выбраны случайно. Каждое сочетание размера и доли дефолтов повторяется с тремя значениями начального состояния генератора, что даёт 60 обучающих подвыборок и позволяет оценить разброс результатов.

Набор содержит только заявки, по которым банк выдал кредит, клиентов с отказом в данных нет. Поэтому ошибка II рода (отказ платёжеспособному клиенту) оценивается среди заявок, одобренных банком, а не среди всех обратившихся, а выборка в целом смещена. Признаки обезличены, что ограничивает содержательную интерпретацию.

\subsection{Метрики и протокол}

Положительным классом считается дефолт. Модель выдаёт вероятность дефолта $p$; заявке отказывают, если $p$ не меньше порога~$t$. Ошибка I рода~--- кредит выдан клиенту, который ушёл в дефолт (главная ошибка), ошибка II рода~--- отказ платёжеспособному клиенту:
\begin{equation}
E_{\mathrm{I}}=\frac{FN}{TP+FN},\qquad
E_{\mathrm{II}}=\frac{FP}{FP+TN},
\end{equation}
где $FN$~--- число дефолтов, которым кредит выдан, $FP$~--- число платёжеспособных клиентов, которым отказано. Обе ошибки выражены долями от объёма своего класса. Дополнительно используются площадь под ROC-кривой (ROC AUC), не зависящая от порога, и площадь под кривой точность--полнота (PR-AUC).

Вероятности моделей, обученных на выборках с разной долей дефолтов, смещены относительно друг друга, поэтому единый порог вида $0{,}5$ для сравнения непригоден. Для всех моделей порог выбирается по одному правилу: на валидации находится порог $t_q$, при котором отказывают доле $q$ наиболее рискованных заявок, $q\in\{5\%;\ 10\%;\ 20\%\}$:
\begin{equation}
t_q=\mathrm{Quantile}_{1-q}\bigl(p_{\mathrm{val}}\bigr).
\end{equation}
Этот порог применяется на тесте. Из-за временного дрейфа реальная доля отказов на тесте может отличаться от $q$, поэтому она фиксируется для каждого запуска.

Параметры моделей не подбираются. Для CatBoost, LightGBM и XGBoost во всех запусках использованы одинаковые значения: 500 деревьев, шаг обучения 0,05, максимальная глубина 6; ранняя остановка не применяется, чтобы валидация не участвовала в обучении моделей, так как на ней подбираются пороги менеджера. Категориальные признаки передаются средствами самих библиотек. TabPFN-2.5 \cite{hollmann2025} использована с параметрами по умолчанию, версия модели задана явно, так как в текущих версиях пакета (\texttt{tabpfn}~9.2.0) по умолчанию загружается более новая модель; расчёт выполнен на графическом ускорителе NVIDIA Tesla~T4 (15\,ГБ), лицензия модели некоммерческая.

Вывод TabPFN на 50\,000 объектов валидации и теста для каждого запуска слишком дорог, поэтому она предсказывает вероятности на фиксированных вложенных случайных подвыборках валидации и теста по [10\,000] объектов. Для сопоставимости метрики бустингов пересчитываются на тех же объектах по сохранённым вероятностям. Различия между моделями оцениваются парным бутстрепом по объектам теста (95\%-й доверительный интервал разности ROC AUC) совместно с разбросом по трём начальным состояниям обучающих подвыборок; разность считается значимой, если интервал не содержит нуль. Для каждого запуска также фиксируются время обучения и время предсказания. Бустинги считались на центральном процессоре, TabPFN~--- на графическом ускорителе, поэтому время разных семейств моделей напрямую не сопоставимо.

% ============================================================
\section{Результаты}
\label{sec:results}
% TODO: профили ошибок моделей, найденные пороги, сравнение с лучшей одиночной моделью, усреднением и стекингом.

% ============================================================
\section{Заключение}
\label{sec:concl}
% TODO: выводы и ограничения (смещение выборки, обезличенные признаки).

% ============================================================
\begin{thebibliography}{9}

\bibitem{xgboost} Chen T., Guestrin C. XGBoost: A scalable tree boosting system. In: \textit{Proc. 22nd ACM SIGKDD Int. Conf. on Knowledge Discovery and Data Mining}, 2016, pp.~785--794.

\bibitem{lightgbm} Ke G. et al. LightGBM: A highly efficient gradient boosting decision tree. In: \textit{Advances in Neural Information Processing Systems}, vol.~30, 2017.

\bibitem{catboost} Prokhorenkova L. et al. CatBoost: unbiased boosting with categorical features. In: \textit{Advances in Neural Information Processing Systems}, vol.~31, 2018.

\bibitem{hollmann2025} Hollmann N. et al. Accurate predictions on small data with a tabular foundation model. \textit{Nature}, vol.~637, pp.~319--326, 2025. DOI: 10.1038/s41586-024-08328-6.

\bibitem{baesens2026} Baesens B. et al. Foundation Models for Credit Risk Prediction: A Game Changer? \textit{arXiv:2605.18147}, 2026. DOI: https://doi.org/10.48550/arXiv.2605.18147

\bibitem{wang2024} Wang, S.; Zhang, X. Research on Credit Default Prediction Model Based on TabNet-Stacking. Entropy 2024, 26, 861. DOI: https://doi.org/10.3390/e26100861

\bibitem{repo} Репозиторий проекта. URL: https://github.com/colziks/credit-scoring-model-manager

\end{thebibliography}

\end{document}
